% STATUS: written 2026-11-05 (Bresar--Chebotar--Martindale 2007 read at
% its preface and its first d-free pages; excerpts in excerpts/. The book
% body is scanned, so only the OCR-transcribed pages are reported. The
% deeper constructive chapters are named only through the Preface.)
\section{d-free sets and applications}\label{sec:functional-identities:d-free}

The $d$-free set of
Section~\ref{sec:functional-identities:fi-theory} is the organizing idea
of the whole theory: it is the basic concept upon which everything is
based \cite[p.~x]{BresarMartindaleChebotar2007FI}. A ring is $d$-free when
two designated two-variable functional identities of degree $d$ admit only
the standard solution \cite[p.~21]{BresarMartindaleChebotar2007FI}. The
concept is introduced exactly for this purpose, to control when
nonstandard solutions can exist
\cite[p.~21]{BresarMartindaleChebotar2007FI}.

The development that the sources report is threefold
\cite[p.~x]{BresarMartindaleChebotar2007FI}. First, the basic properties
and constructions of $d$-free sets are worked out
\cite[p.~x]{BresarMartindaleChebotar2007FI}. Second, functional
identities on \emph{d}-free sets are analyzed in general, through the
notions of quasi-polynomials and core functions
\cite[p.~x]{BresarMartindaleChebotar2007FI}. Third, the existence of
central classes of $d$-free sets is proved inside prime and semiprime
rings, which are the rings the applications live on
\cite[p.~x]{BresarMartindaleChebotar2007FI}.

What the concept buys is a single hypothesis that replaces the separate
structural arguments of
Section~\ref{sec:functional-identities:commuting-maps}. The commuting map
$f(x)=\lambda x+\mu(x)$ of Theorem~\ref{thm:fi:leetheorem2} was obtained
ring by ring. On a sufficiently $d$-free ring the same conclusion follows
from the general theory of functional identities on $d$-free sets
\cite[p.~x]{BresarMartindaleChebotar2007FI}. The applications reported in
the sources are of two families: to derivation problems, taken up in
Section~\ref{sec:functional-identities:applications-to-derivations}, and
to centralizers, multipliers and Lie centralizers, taken up in
Section~\ref{sec:functional-identities:centralizers-multipliers}
\cite[p.~x]{BresarMartindaleChebotar2007FI}.
